\ssszh 转到图 \ref{f:tableaut} 与图 \ref{f:tablespec} 的第三行，研究以 Whittaker 模型为源的缠绕算子。先遇到它的自同态，以及一个幺半范畴的等价
%
\begin{equation} \label{e:whitend} \Dmod(LN, \psi \backslash LG / LN, \psi)
\simeq \QCoh(\LS).\end{equation}
%
这是一个引人注目、且基本上仍完全开放的猜想。甚至连左端的对称幺半结构都远非显然。从精神上说，
在当前几何设定下，它是 Whittaker 余向量唯一性的一种体现——在通常的「重数自由」与「自同态代数交换」的通常对应之下。

我们强调，这一等价会确定猜想 \eqref{e:loclang} 中一大片内容，即谱侧所有奇异支撑沿零截面的
范畴层。这样的范畴表示称为缓增的。

\begin{re} 尽管 \eqref{e:whitend} 基本上仍是开放的，这里有一个已知的特例。也就是说，由 Raskin \cite{whit}
的工作，环路群的范畴表示上的深度滤过给出
\[
\Dmod(LN, \psi \backslash LG / LN, \psi)
\]
上的一个滤过，其中每个商都是幺半余局部化。第一个商对应「向量被 Iwahori 子群的
射幂幺根固定的表示」，即 Iwahori--Whittaker 范畴
%
\[
     \Dmod(\mathring{I}, \psi \backslash LG / \mathring{I}, \psi).
\]
%
它又容许一个进一步的商，后者对应从全部驯分歧表示过渡到幂幺单值的那些表示，即
Steinberg--Whittaker 范畴
%
\[
   \Dmod(\mathring{I}, \psi, \on{s} \backslash LG / \mathring{I}, \psi,
   \on{s}).
\]
%
显然，它就是由 $\Dmod(I^- \backslash LG / I^-)$ 中对象的平均所生成的满子范畴，其中 $I^-$ 是
$L^+G$ 中相对于与 $B$ 处于一般位置的 Borel 子群 $B^-$ 定义的 Iwahori 子群，参见 \cite{ahc}。
若以 $\sN^{\wedge}$ 记 $\cfg$ 沿幂幺锥 $\sN$ 的形式完备化，则有一个幺半等价
%
\[
    \Dmod(\mathring{I}, \psi, \on{s} \backslash LG / \mathring{I}, \psi,
    \on{s}) \simeq \QCoh( \sN^\wedge / \cG),
\]
%
这正是 \eqref{e:whitend} 所许诺的一块。

Bezrukavnikov 证明仿射 Hecke 范畴的等价 \eqref{e:2realz}，本质上就是把这一结果与前面提到的
Arkhipov--Bezrukavnikov 结果结合起来。更准确地说，以 $\widetilde{\sN}^\wedge$ 记
Grothendieck--Springer 簇
\[
\widetilde{\cG} := \cG \overset \cB \times \check{\fb}
\]
沿 Springer 簇的形式完备化。也就是说，Grothendieck--Springer 簇参数化数对 $(X, \check{\fb}')$，
其中 $X$ 是 $\cfg$ 的元素，$\check{\fb}'$ 是包含它的 Borel 子代数；而沿 $X$ 幂幺的轨迹作完备化。
于是所要的等价 \eqref{e:2realz}\footnote{更仔细地说，我们指的是它的单值变体；严格的陈述可以从它
形式地推出。}由重整化下述等价得到
%
{\small
\begin{align*} &\QCoh( \widetilde{\sN}^\wedge / \cG  \underset {\cfg / \cG}
\times
\widetilde{\sN}^\wedge / \cG) \\
&\qquad \simeq \QCoh(\widetilde{\sN}^\wedge / \cG)  \underset {\QCoh(\sN^\wedge / \cG)}
\otimes  \QCoh(\widetilde{\sN}^\wedge / \cG ) \\ &\qquad \simeq \Dmod(\mathring{I}
\backslash LG / \mathring{I}, \psi, \on{s}) \underset {\Dmod(\mathring{I},
\psi, \on{s} \backslash LG / \mathring{I}, \psi, \on{s})} \otimes
\Dmod(\mathring{I}, \psi, \on{s} \backslash LG / \mathring{I}).  \end{align*}}
%

\end{re}

\ssszh 接下来把 Whittaker 模型与 Kac--Moody 表示配对，得到等价
\begin{equation} \label{e:affskr} \fgc\mod^{LN, \psi} \simeq \QCoh(\Op).
\end{equation}
这同样是一条关于任意分歧对象的强断言。与我们正遇到的其他这类断言不同，它事实上是 Raskin
\cite{whit} 的一条定理。

也就是说，Raskin 的仿射 Skryabin 定理\footnote{附带一提，Skryabin 是「by Raskin」的变位词。}
对任意水平 $\kappa$ 给出一个典范等价
%
\[
  \fgk\mod^{LN, \psi} \simeq \sW\mod,
\]
%
其中 $\sW$ 是仿射 W-代数。另一方面，在临界水平，Feigin--Frenkel 已把 $\mathscr{W}_{\kappa_c}$ 等同于
包络代数的中心 $Z(\fgc)$ \cite{feiginfrenkel91}，因而也等同于 $\OO_{\Op}$。这与仿射 Skryabin 定理
结合，就给出 \eqref{e:affskr}。

\begin{re} 读者可参考本文集中 Arakawa 的综述，以了解 W-代数的更多细节；另见 \cite{araintro}、\cite{fbz}。
这里只附带一提：仿射 Skryabin 定理使得人们可以在 W-代数表示的研究中系统地使用范畴表示论的
技巧。例如可参见 \cite{locw}，其中讨论了 W-代数的最高权模的局部化理论。
\end{re}

\ssszh \label{ss:affhc}最后一行遇到的是 Kac--Moody 表示的自同态。它们如下等同于临界水平处的
仿射 Harish-Chandra 双模。

首先，范畴 $\fgc\mod$ 作为 $\Dmod(LG)$ 表示、特别地作为抽象 dg 范畴，典范地自对偶。

为看清这一点，从更简单的情形入手会有帮助。对带中心扩张的有限维李代数 $\mathfrak{a}$
%
\[
0 \rightarrow k \cdot \mathbf{1} \rightarrow \widetilde{\mathfrak{a}}_c
\rightarrow \mathfrak{a} \rightarrow 0,
\]
%
考虑 $\mathbf{1}$ 通过恒等作用的 $\widetilde{\mathfrak{a}}_c$-模范畴 $\widetilde{\mathfrak{a}}_c\mod$。
若以 $\widetilde{\mathfrak{a}}_{-c}$ 记相反的中心扩张，例如由任一表示 2-余循环的加法逆给出，则有一个
完美配对
%
\[
\widetilde{\mathfrak{a}}_c\mod \otimes \hspace{.7mm}
\widetilde{\mathfrak{a}}_{-c}\mod \rightarrow \mathfrak{a}\mod \rightarrow
\Vect,
\]
%
其中第一个箭头把表示张量起来以消去扩张，第二个箭头是李代数同调。

在仿射情形，Arkhipov--Gaitsgory \cite{ag} 在 Frenkel--Gaitsgory \cite{fg06loc} 先前工作的基础上证明了：
对任意 $\kappa$，复合
%
\[
\fgk\mod \otimes \hspace{.7mm} \widehat{\fg}_{-\kappa + 2\kappa_c}\mod
\rightarrow \widehat{\fg}_{2 \kappa_c}\mod \rightarrow \Vect
\]
%
是完美配对，其中第一个箭头把表示张量起来，第二个箭头是半无限上同调函子。简言之，后者是像 $\gk$
这样的无限维李代数（正式地说，Tate 李代数）的表示的自然同调理论，但它需要水平为 $2\kappa_c$ 的中心
扩张。对一般 Tate 李代数也有这样的对偶陈述，参见 \cite{mys}、\cite{lpw}。

特别地，有自对偶
%
\[
    \fgc\mod \otimes \hspace{.7mm} \fgc\mod \rightarrow \Vect,
\]
%
这是「临界水平是水平空间的中点」的另一种意义。此外，如第 \ref{ss:hchbimod} 节那样，这把
$\fgc\mod$ 的 $\Dmod(LG)$-等变自函子等同于临界水平仿射 Harish-Chandra 双模范畴
%
\[
      \fgc \oplus \fgc \mod^{LG}.
\]
%
这是一个漂亮却也多少难以把握的范畴。例如，它的全部对象相对于 $\fgc \oplus \fgc\mod$ 上自然的
$t$-结构都落在同调次数 $-\infty$，也就是说，它由通过重整化引入的同调幽灵构成。

\begin{re} 其他水平处的仿射 Harish-Chandra 双模同样微妙，其中一些基础工作与猜想属于 I. Frenkel--Malikov
\cite{frenkelmalikov1} 与 Gaitsgory \cite{quantum-langlands-summary}。
\end{re}

对这些双模，局部几何朗兰兹预测一个与谱卷积代数的幺半等价
%
\[
      \fgc \oplus \fgc \mod^{LG} \simeq \QCoh( \Op \underset \LS \times \Op).
\]
%
这是一个引人注目、且基本上仍完全开放的猜想。据我们所知，甚至连球情形与幂幺单值的驯分歧情形都还是
开放的；不过我们预期，利用前面提到的 Frenkel--Gaitsgory 的一些工作（参见第 \ref{ss:sphericalkmmod}
与 \ref{ss:iwahorikmmod} 节），再结合临界水平的局部化定理，这两个情形是可以达到的。下面转入该定理。
